\par\Needspace{6\baselineskip}\medskip
\begin{center}\hypertarget{native-189-article141}{}\pdfbookmark[2]{141.}{article-141}141.\end{center}
From $P=\dfrac{n''}{n}$ and equation III of the preceding article it follows that
\[
(n+n'')\frac{P+a}{P+1}=bn'\frac{\sin(z-\sigma)}{\sin z}\mathpunct{\raisebox{.45ex}{;}}
\]
From this and from $Q=2\left(\dfrac{n+n''}{n'}-1\right)r'^3$ and $r'=\dfrac{R'\sin\delta'}{\sin z}$ we obtain
\[
\sin z+\frac{Q\sin z^4}{2R'^3\sin\delta'^3}=b\frac{P+1}{P+a}\sin(z-\sigma)
\]
or
\[
\begin{aligned}
\frac{Q\sin z^4}{2R'^3\sin\delta'^3}&=\left(b\frac{P+1}{P+a}-\cos\sigma\right)\sin(z-\sigma)\\[.5ex]
&\quad-\sin\sigma\cos(z-\sigma).
\end{aligned}
\]
Hence, if for brevity we set
\[
\tag*{14)}\frac{1}{2R'^3\sin\delta'^3\sin\sigma}=c
\]
and introduce an auxiliary angle $\omega$ such that
\[
\operatorname{tang}\omega=\frac{\sin\sigma}{b\dfrac{P+1}{P+a}-\cos\sigma}\mathpunct{\raisebox{.45ex}{,}}
\]
we obtain the equation
\[
\tag*{IV.}cQ\sin\omega\sin z^4=\sin(z-\omega-\sigma),
\]
\SourceLeaf{190}from which the unknown $z$ must be determined. To compute the angle $\omega$ more easily, it is convenient to write the preceding formula for $\operatorname{tang}\omega$ in the form
\[
\operatorname{tang}\omega=\frac{(P+a)\operatorname{tang}\sigma}{P\left(\dfrac{b}{\cos\sigma}-1\right)+\left(\dfrac{b}{\cos\sigma}-a\right)}\mathpunct{\cdot}
\]
Therefore, on setting
\[
\tag*{15)}\frac{\dfrac{b}{\cos\sigma}-a}{\dfrac{b}{\cos\sigma}-1}=d
\]
\[
\tag*{16)}\frac{\operatorname{tang}\sigma}{\dfrac{b}{\cos\sigma}-1}=e
\]
we obtain the following very simple formula for determining $\omega$
\[
\operatorname{tang}\omega=\frac{e(P+a)}{P+d}\mathpunct{\cdot}
\]
We regard the calculation of the quantities $a,b,c,d,e$ by formulas 11–16, which depends only on the given quantities, as the fourth task. For $b,c,e$, only their logarithms will be needed.

There is a special case in which these rules require some modification. Whenever the great circle $BB''$ coincides with $A''B''$, so that the points $B,B^\ast$ coincide respectively with $D',D$, the quantities $a,b$ would take infinite values. In this case, if we set
\[
\frac{R\sin\delta\sin(A'D''-\delta'+\sigma)}{R'\sin\delta'\sin(AD''-\delta)}=\pi,
\]
we obtain, in place of equation III, $0=\pi n-\dfrac{n'\sin(z-\sigma)}{\sin z}\mathpunct{\raisebox{.45ex}{,}}$ and by setting
\[
\operatorname{tang}\omega=\frac{\pi\sin\sigma}{P+(1-\pi\cos\sigma)}
\]
we obtain the same equation IV.

Similarly, in the special case where $\sigma=0$, $c$ becomes infinite and $\omega=0$, so that the factor $c\sin\omega$ in equation IV appears indeterminate. It is in fact determinate, and its value is
\[
=\frac{P+a}{2R'^3\sin\delta'^3(b-1)(P+d)}\mathpunct{\raisebox{.45ex}{,}}
\]
as a moment's consideration shows. Thus in this case
\[
\sin z=R'\sin\delta'\;\BareRoot[3]{\frac{2(b-1)(P+d)}{Q(P+a)}}\mathpunct{\cdot}
\]
